%% GENERATED by tools/gen_latex.py -- edit content/ch12.py instead.
%% Build:  latexmk -pdf ch12.tex     (or: pdflatex ch12.tex, twice)
\documentclass[aspectratio=169,11pt,t]{beamer}
\usepackage{itbpro}

\coursetitle{Object Detection}
\coursesubtitle{Boxes and labels, predicted together --- a simplified YOLO built from scratch, and a pretrained detector that generalises to a pointillist painting.}
\coursekicker{Chapter 12}
\coursesource{Chollet \& Watson, Deep Learning with Python 3e --- chapter 12}
\courseduration{2.5 hours}
\coursepresenter{Rahman Indra Kesuma, S.Kom., M.Cs.}{Teaching Assistant}
\coursebrand{AI for Professional \textperiodcentered\ ITB \textperiodcentered\ Ch 12}

\title{Object Detection}
\author{Rahman Indra Kesuma, S.Kom., M.Cs.}
\date{}

\begin{document}
\covertitle

\framekicker{Learning outcomes}
\begin{frame}[shrink=25]{Session Objectives}
\leadpar{By the end of this session, participants can:}
\begin{isteps}
\item Say when to use \textbf{detection rather than segmentation}, on both compute and labelling grounds.
\item Distinguish \textbf{two-stage (R-CNN)} from \textbf{single-stage (YOLO, RetinaNet, SSD)} detectors, and name the trade-off.
\item Build a \textbf{YOLO detection head} on a pretrained backbone, and read the shape of its grid output.
\item Explain why the \textbf{confidence target is an IoU score} rather than a 0/1 flag.
\item Run a \textbf{pretrained RetinaNet}, and explain why it handles small and large objects better.
\item Say why single-stage detectors are \textbf{more robust to unfamiliar inputs} than two-stage ones.
\end{isteps}

\vskip 6pt
\vskip 4pt
\reslink{Book}{Chapter 12 --- full text}{https://deeplearningwithpython.io/chapters/chapter12\_object-detection}
\reslink{Notebook}{01\_coco\_and\_boxes.ipynb}{../../notebooks/ch12/01\_coco\_and\_boxes.ipynb}
\reslink{Notebook}{02\_yolo\_from\_scratch.ipynb}{../../notebooks/ch12/02\_yolo\_from\_scratch.ipynb}
\reslink{Notebook}{03\_pretrained\_retinanet.ipynb}{../../notebooks/ch12/03\_pretrained\_retinanet.ipynb}
\reslink{Repo}{Author's official notebook}{https://github.com/fchollet/deep-learning-with-python-notebooks/blob/master/chapter12\_object-detection.ipynb}
\vskip 2pt
\end{frame}

\framekicker{Introduction}
\begin{frame}[shrink=25]{What detection is for}
{\fontsize{9.4}{12.6}\selectfont\color{fg2}A detector draws boxes around objects and labels them, so you know not just \textbf{which} objects are present but \textbf{where}.\par}\vskip 4pt
\begin{minipage}[t]{0.3253\textwidth}
\begin{icard}
\cardh{Counting}
How many instances of an object are in this image.
\end{icard}
\end{minipage}%
\hfill
\begin{minipage}[t]{0.3253\textwidth}
\begin{icard}
\cardh{Tracking}
Run detection on every frame to follow objects over time.
\end{icard}
\end{minipage}%
\hfill
\begin{minipage}[t]{0.3253\textwidth}
\begin{icard}
\cardh{Cropping}
Cut out the region containing an object and send a higher-resolution patch to a classifier or an OCR model.
\end{icard}
\end{minipage}%
\end{frame}

\framekicker{Introduction}
\begin{frame}[shrink=25]{If segmentation is a superset, why detect at all?}
{\fontsize{9.4}{12.6}\selectfont\color{fg2}A fair question: given an instance mask you can already compute the smallest box containing it. \textbf{Segmentation is a strict superset of detection.}\par}\vskip 4pt
\begin{minipage}[t]{0.494\textwidth}
\begin{icardgood}
\cardh{Computational cost}
A good detector typically runs \textbf{much faster} than a segmentation model.
\end{icardgood}
\end{minipage}%
\hfill
\begin{minipage}[t]{0.494\textwidth}
\begin{icardgood}
\cardh{Labelling cost}
Segmentation needs \textbf{pixel-precise masks}, far more time-consuming to produce than bounding boxes.
\end{icardgood}
\end{minipage}%
\begin{iband}[signal]
So the rule: \textbf{always use a detector if you do not need pixel-level information} --- for instance if all you want is \hilite{to count things}
\end{iband}
\end{frame}

\coursesection{01}{Two-stage vs single-stage}{Two broad families, and why one of them won.}

\framekicker{Section 12.1.1}
\begin{frame}[shrink=25]{Two-stage: propose, then classify}
\begin{center}
\includegraphics[width=0.92\linewidth,height=0.60\textheight,keepaspectratio]{figs/ch12-twostage.tex.pdf}
\end{center}
\vskip 3pt{\color{ink3}\fontsize{7.4}{9.4}\selectfont Figure 12.2 --- an R-CNN extracts region proposals, then classifies each one with a ConvNet.\par}
{\fontsize{9.4}{12.6}\selectfont\color{fg2}The first stage is \textbf{not very smart}: it produces boxes around areas that merely \textit{look object-like}. The second stage decides what, if anything, is in each one.\par}\vskip 4pt
\end{frame}

\framekicker{Section 12.1.1}
\begin{frame}[shrink=25]{Why that is expensive}
\begin{iband}[amber]
It works well, but it requires classifying \textbf{thousands of patches for every single image}. That rules it out for most real-time applications and for embedded systems.
\end{iband}
{\fontsize{9.4}{12.6}\selectfont\color{fg2}The book's practical take is blunt: you rarely need R-CNN at all. \textbf{With a large server GPU you are usually better off with a segmentation model} like SAM; \textbf{if you are resource-constrained you want a single-stage detector}. \hilite{Neither case points at a two-stage model.}\par}\vskip 4pt
\end{frame}

\framekicker{Section 12.1.2}
\begin{frame}[shrink=25]{Single-stage: predict both at once}
\begin{center}
\includegraphics[width=0.92\linewidth,height=0.60\textheight,keepaspectratio]{figs/ch12-singlestage.tex.pdf}
\end{center}
\vskip 3pt{\color{ink3}\fontsize{7.4}{9.4}\selectfont One model jointly predicts box coordinates and labels.\par}
{\fontsize{9.4}{12.6}\selectfont\color{fg2}The main families are \textbf{RetinaNet}, \textbf{SSD} (Single Shot MultiBox Detector), and the \textbf{YOLO} family --- \textit{You Only Look Once}, named after the meme on purpose.\par}\vskip 4pt
\end{frame}

\framekicker{Section 12.1.2}
\begin{frame}[shrink=25]{The trade-off, and the state of play}
\begin{ibullets}
\item Single-stage detectors are \textbf{significantly faster and more efficient} than two-stage ones, with a \textbf{minor potential trade-off in accuracy}.
\item \textbf{YOLO is arguably the most popular object detection model there is}, especially for real-time applications.
\item A new version appears roughly every year --- and interestingly, \hilite{each new version tends to come from a different organisation}.
\end{ibullets}
{\fontsize{9.4}{12.6}\selectfont\color{fg2}There were twelve YOLO versions at the time of writing. This chapter recreates \textbf{the original from 2015}, which is simpler to work with.\par}\vskip 4pt
\end{frame}

\coursesection{02}{Training a YOLO model from scratch}{COCO, a grid, and a loss with an unusual target.}

\framekicker{Section 12.2}
\begin{frame}[shrink=25]{An honest warning before starting}
\begin{iband}[signal]
Building a detector is \textbf{an undertaking} --- not because anything in it is theoretically complex, but because there is \hilite{a lot of code just to manipulate bounding boxes and predicted output}
\end{iband}
{\fontsize{9.4}{12.6}\selectfont\color{fg2}The dataset is \textbf{COCO} --- \textit{Common Objects in Context} --- one of the best-known detection datasets. It ships object labels, bounding boxes, and full segmentation masks; this chapter uses only the boxes.\par}\vskip 4pt
\end{frame}

\framekicker{Section 12.2.1 $\cdot$ listing 12.1}
\begin{frame}[shrink=25]{Downloading COCO}
{\fontsize{9.4}{12.6}\selectfont\color{fg2}Two archives: the images, and a JSON of annotations. At \textbf{18 GB} this is the largest dataset in the book --- though the book notes it is not large by modern standards.\par}\vskip 4pt
\icode{python}{listing 12.1 --- the 2017 COCO dataset}{listings/ch12-01.py}
{\fontsize{9.4}{12.6}\selectfont\color{fg2}The first gives an \textbf{unlabelled directory of images}; the second gives all the metadata. They have to be joined.\par}\vskip 4pt
\end{frame}

\framekicker{Section 12.2.1}
\begin{frame}[shrink=25]{Joining images to their boxes}
\begin{center}
\includegraphics[width=0.92\linewidth,height=0.60\textheight,keepaspectratio]{figs/ch12-coco.tex.pdf}
\end{center}
\vskip 3pt{\color{ink3}\fontsize{7.4}{9.4}\selectfont COCO associates each image file with an ID, and each bounding box with one of those IDs.\par}
\end{frame}

\framekicker{Section 12.2.1 $\cdot$ listing 12.2}
\begin{frame}[shrink=25]{\ldots{}and normalising the coordinates while you are at it}
{\fontsize{9.4}{12.6}\selectfont\color{fg2}Boxes arrive as pixel coordinates from the top-left corner. Rescaling them into a unit square means the rest of the code \textbf{never has to check an image's size}.\par}\vskip 4pt
\icode{python}{listing 12.2 --- parsing the annotations}{listings/ch12-02.py}
\begin{iband}[signal]
The two centring lines matter: images are not square, so scaling by the longest side leaves letterboxing --- and the boxes have to be \hilite{shifted by half that gap} or every prediction is offset.
\end{iband}
\end{frame}

\framekicker{Section 12.2.2}
\begin{frame}[shrink=25]{The core idea: a grid of small predictors}
\begin{center}
\includegraphics[width=0.92\linewidth,height=0.60\textheight,keepaspectratio]{figs/ch12-yologrid.tex.pdf}
\end{center}
\vskip 3pt{\color{ink3}\fontsize{7.4}{9.4}\selectfont Figure 12.4 --- each grid cell emits a box, a confidence, and a class distribution.\par}
{\fontsize{9.4}{12.6}\selectfont\color{fg2}The original paper predicted several boxes per cell; this version keeps \textbf{one box per cell} for simplicity.\par}\vskip 4pt
\end{frame}

\framekicker{Section 12.2.2}
\begin{frame}[shrink=25]{Why a confidence score is needed at all}
{\fontsize{9.4}{12.6}\selectfont\color{fg2}Objects are not evenly distributed across a grid. \textbf{Most cells contain nothing}, and the model needs a way to say so.\par}\vskip 4pt
\begin{iband}[signal]
So each cell also outputs a \textbf{confidence}: high where an object is detected, near zero where there is none. \hilite{Most cells should report near-zero confidence}, and that is the correct behaviour, not a failure.
\end{iband}
\end{frame}

\framekicker{Section 12.2.2 $\cdot$ listing 12.5}
\begin{frame}[shrink=25]{The backbone: ResNet, not Xception}
{\fontsize{9.4}{12.6}\selectfont\color{fg2}Same KerasHub pattern as chapter 8, with two deliberate changes: a different architecture, and a larger input size.\par}\vskip 4pt
\icode{python}{listing 12.5 --- loading ResNet}{listings/ch12-03.py}
\begin{ibullets}
\item \textbf{ResNet} is structurally similar to Xception but downsamples with \textbf{strides instead of pooling} --- and chapter 11 established that strides preserve location, which detection needs.
\item \textbf{448 $\times$ 448} rather than 180 $\times$ 180: \hilite{input size matters a great deal} for detection, because small objects vanish at low resolution.
\end{ibullets}
\end{frame}

\framekicker{Section 12.2.2 $\cdot$ listing 12.6}
\begin{frame}[shrink=25]{The detection head}
{\fontsize{9.4}{12.6}\selectfont\color{fg2}The head proposed in the YOLO paper is simple: take the backbone output, push it through two dense layers, then split the result.\par}\vskip 4pt
\icode{python}{listing 12.6 --- attaching the head}{listings/ch12-04.py}
{\fontsize{9.4}{12.6}\selectfont\color{fg2}The \inlinecode{+ 5} is four box coordinates plus one confidence. Everything after that is the class distribution.\par}\vskip 4pt
\end{frame}

\framekicker{Section 12.2.2}
\begin{frame}[shrink=25]{Following the shapes through the head}
\begin{center}
\includegraphics[width=0.92\linewidth,height=0.60\textheight,keepaspectratio]{figs/ch12-yolohead.tex.pdf}
\end{center}
\vskip 3pt{\color{ink3}\fontsize{7.4}{9.4}\selectfont The strided convolution exists purely to get the float count down to something a dense layer can accept.\par}
\begin{minipage}[t]{0.494\textwidth}
\istat{401,408}{floats per image out of the backbone}
\end{minipage}%
\hfill
\begin{minipage}[t]{0.494\textwidth}
\istat{18,432}{after the strided convolution}
\end{minipage}%
\end{frame}

\framekicker{Section 12.2.2}
\begin{frame}[shrink=25]{One design choice worth noticing}
\begin{iband}[signal]
Because the \textbf{entire feature map is flattened} before the dense layers, every grid detector can see features from \textbf{the whole image}. There is \hilite{no locality constraint}
\end{iband}
{\fontsize{9.4}{12.6}\selectfont\color{fg2}That is deliberate: \textbf{large objects will not stay contained inside a single grid cell}, so a cell responsible for one needs to see beyond its own square.\par}\vskip 4pt
\begin{minipage}[t]{0.494\textwidth}
\istat{77.8 M}{total parameters}
\end{minipage}%
\hfill
\begin{minipage}[t]{0.494\textwidth}
\istat{297 MB}{model size}
\end{minipage}%
\end{frame}

\framekicker{Section 12.2.3}
\begin{frame}[shrink=25]{Aligning the labels with the grid}
{\fontsize{9.4}{12.6}\selectfont\color{fg2}The model is simple; the data preparation is where the work is. Targets have to be rearranged into the same 6 $\times$ 6 grid the model predicts.\par}\vskip 4pt
\begin{iband}[signal]
The assignment rule: \textbf{each grid detector is responsible for any box whose centre falls inside its cell}. A box straddling several cells still belongs to \hilite{exactly one of them}
\end{iband}
{\fontsize{9.4}{12.6}\selectfont\color{fg2}This is why the box format matters: the coordinates are stored as centre-x, centre-y, width, height, so the cell assignment is a direct lookup.\par}\vskip 4pt
\end{frame}

\framekicker{Section 12.2.3}
\begin{frame}[shrink=25]{Two coordinate systems, and the pair of functions between them}
\begin{center}
\includegraphics[width=0.92\linewidth,height=0.60\textheight,keepaspectratio]{figs/ch12-gridcoord.tex.pdf}
\end{center}
\vskip 3pt{\color{ink3}\fontsize{7.4}{9.4}\selectfont The model predicts positions relative to a cell; the labels arrive relative to the image.\par}
{\fontsize{9.4}{12.6}\selectfont\color{fg2}\textbf{\inlinecode{x} and \inlinecode{y} are relative to the grid cell} (0 to 1 within it), while \textbf{\inlinecode{w} and \inlinecode{h} stay relative to the whole image}. The widths and heights need no conversion; \hilite{only the centres do}\par}\vskip 4pt
\end{frame}

\framekicker{Section 12.2.3}
\begin{frame}[shrink=25]{\ldots{}written out}
{\fontsize{9.4}{12.6}\selectfont\color{fg2}Two small functions, exact inverses of one another --- one to prepare targets, one to read predictions back.\par}\vskip 4pt
\icode{python}{converting to and from the grid}{listings/ch12-05.py}
\begin{iband}[amber]
Getting one of these subtly wrong is the classic detection bug: the loss still falls, the boxes still appear, and they are \hilite{consistently in the wrong place}
\end{iband}
\end{frame}

\framekicker{Section 12.2.3 $\cdot$ listing 12.7}
\begin{frame}[shrink=25]{Building the two target arrays}
{\fontsize{9.4}{12.6}\selectfont\color{fg2}One array holds the class map, the other the boxes with their confidence --- and the confidence in the \textbf{labels} is always exactly 1 or 0.\par}\vskip 4pt
\icode{python}{listing 12.7 --- creating the targets}{listings/ch12-06.py}
\begin{iband}[signal]
Note the asymmetry: \textbf{the class map marks every cell the box overlaps}, while \textbf{the box array marks only the cell containing its centre}. Overlapping boxes are \hilite{deliberately not handled}, to keep the example readable.
\end{iband}
\end{frame}

\framekicker{Section 12.2.4}
\begin{frame}[shrink=25]{Two losses, summed}
{\fontsize{9.4}{12.6}\selectfont\color{fg2}The model has two outputs, so it needs two losses --- exactly the multi-output pattern from chapter 7, where \textbf{Keras sums them during training}.\par}\vskip 4pt
\begin{itable}{>{\raggedright\arraybackslash}p{0.188\textwidth}>{\raggedright\arraybackslash}p{0.3196\textwidth}>{\raggedright\arraybackslash}p{0.4324\textwidth}}
\thead{Output} & \thead{Loss} & \thead{Note} \\ \midrule
\textbf{class} & \inlinecode{sparse\_categorical\_crossentropy} & Entirely ordinary. \\
\textbf{box} & Sum-squared error, custom & Computed \textbf{only for grid cells that actually contain a box} in the labels. \\
\end{itable}
\end{frame}

\framekicker{Section 12.2.4}
\begin{frame}[shrink=25]{The clever part: what confidence is trained against}
\begin{center}
\includegraphics[width=0.92\linewidth,height=0.60\textheight,keepaspectratio]{figs/ch12-conf.tex.pdf}
\end{center}
\vskip 3pt{\color{ink3}\fontsize{7.4}{9.4}\selectfont The confidence target is not a flag --- it is a measure of how good the box currently is.\par}
{\fontsize{9.4}{12.6}\selectfont\color{fg2}The authors wanted confidence to reflect \textbf{not just the presence of an object but how good the predicted box is}. So as the model gets better at locating boxes, \hilite{the IoU rises and the learned confidence rises with it}\par}\vskip 4pt
\end{frame}

\framekicker{Section 12.2.4 $\cdot$ listing 12.10}
\begin{frame}[shrink=25]{Computing IoU for a batch of boxes}
{\fontsize{9.4}{12.6}\selectfont\color{fg2}Chapter 11 used IoU as a metric; here it has to be differentiable and run inside the loss, so it is written with \inlinecode{keras.ops}.\par}\vskip 4pt
\icode{python}{listing 12.10 --- intersection area}{listings/ch12-07.py}
\begin{iband}[signal]
The \inlinecode{ops.maximum(0.0, ...)} at the end is what handles \textbf{boxes that do not overlap at all}: a negative width would otherwise produce \hilite{a spurious positive area}
\end{iband}
\end{frame}

\framekicker{Section 12.2.4 $\cdot$ listing 12.13}
\begin{frame}[shrink=25]{Running it, and reading the prediction}
{\fontsize{9.4}{12.6}\selectfont\color{fg2}Predictions come back as a grid, so drawing them means taking an argmax for the label at each cell and filtering by confidence.\par}\vskip 4pt
\icode{python}{listing 12.13 --- visualising a prediction}{listings/ch12-08.py}
{\fontsize{9.4}{12.6}\selectfont\color{fg2}The cutoff is set low because, as the next slide says plainly, \hilite{the model is not a very good detector yet}\par}\vskip 4pt
\end{frame}

\framekicker{Section 12.2.4}
\begin{frame}[shrink=25]{The chapter says so out loud}
\begin{iband}[amber]
Training takes \textbf{over an hour on a free Colab GPU}, and the model is \textbf{still undertrained} --- the validation loss is still falling when it stops. It is starting to understand box locations and labels, \hilite{but it is not accurate}
\end{iband}
{\fontsize{9.4}{12.6}\selectfont\color{fg2}This is worth dwelling on: the book publishes a result that does not work well, and explains why. \textbf{The purpose of the exercise is to feel how detection training behaves}, not to produce a usable COCO detector.\par}\vskip 4pt
\end{frame}
\note{A good teaching moment about honest reporting --- most tutorials would quietly show only the one image that looks convincing.}

\framekicker{Section 12.2.4}
\begin{frame}[shrink=25]{What would actually be needed}
\begin{isteps}
\item \textbf{Train for more epochs} --- the loss had not converged.
\item \textbf{Use the whole COCO dataset}, not the subset.
\item \textbf{Data augmentation} --- translating and rotating both the images \textbf{and their boxes}.
\item \textbf{Improve the class probability map} for overlapping boxes.
\item \textbf{Predict multiple boxes per grid location}, with a larger output grid.
\end{isteps}
\begin{iband}[signal]
Together those approach the original YOLO training recipe. Doing it properly would take \hilite{a large amount of compute and time}, which is the honest reason to reach for a pretrained detector instead.
\end{iband}
\end{frame}

\coursesection{03}{A pretrained RetinaNet}{The same principles, with one important architectural difference.}

\framekicker{Section 12.3}
\begin{frame}[shrink=25]{The difference that matters: scale}
\begin{center}
\includegraphics[width=0.92\linewidth,height=0.60\textheight,keepaspectratio]{figs/ch12-scale.tex.pdf}
\end{center}
\vskip 3pt{\color{ink3}\fontsize{7.4}{9.4}\selectfont RetinaNet uses its ConvNet differently, to handle small and large objects simultaneously.\par}
{\fontsize{9.4}{12.6}\selectfont\color{fg2}In our YOLO we used only the \textbf{final} backbone output. Those features map to large areas of the input, so they are \hilite{not very effective at finding small objects}\par}\vskip 4pt
\end{frame}

\framekicker{Section 12.3}
\begin{frame}[shrink=25]{Loading it, and the box format argument}
{\fontsize{9.4}{12.6}\selectfont\color{fg2}One line to instantiate --- and one argument worth understanding, because box formats are a persistent source of silent errors.\par}\vskip 4pt
\icode{python}{loading RetinaNet}{listings/ch12-09.py}
\begin{iband}[signal]
\inlinecode{rel} means \textbf{relative to the image size} --- coordinates in [0, 1]. Most Keras models and layers that handle boxes accept this argument, and \hilite{setting it wrongly produces boxes in the wrong place with no error at all}
\end{iband}
\end{frame}

\framekicker{Section 12.3}
\begin{frame}[shrink=25]{What comes back}
{\fontsize{9.4}{12.6}\selectfont\color{fg2}Four outputs, structurally similar to the YOLO model's --- with a fixed ceiling of 100 detections per image.\par}\vskip 4pt
\icode{python}{inspecting the prediction}{listings/ch12-10.py}
\ioutput{listings/ch12-11.txt}
{\fontsize{9.4}{12.6}\selectfont\color{fg2}\inlinecode{num\_detections} is the one to read first --- the arrays are padded to 100, so \textbf{the rest is not meaningful}.\par}\vskip 4pt
\end{frame}

\framekicker{Section 12.3 $\cdot$ listing 12.15}
\begin{frame}[shrink=25]{Drawing the detections}
{\fontsize{9.4}{12.6}\selectfont\color{fg2}Iterate only up to \inlinecode{num\_detections}, and translate the numeric COCO label into a name.\par}\vskip 4pt
\icode{python}{listing 12.15 --- running inference}{listings/ch12-12.py}
{\fontsize{9.4}{12.6}\selectfont\color{fg2}That is the whole of using a pretrained detector: load, predict, draw.\par}\vskip 4pt
\end{frame}

\framekicker{Section 12.3}
\begin{frame}[shrink=25]{The pointillist painting test}
\begin{iband}[signal]
RetinaNet generalises to \textbf{a pointillist painting} with ease, despite never having been trained on that style. Paintings and photographs differ enormously at the pixel level but \hilite{share a similar structure at a high level}
\end{iband}
{\fontsize{9.4}{12.6}\selectfont\color{fg2}And the book draws an architectural conclusion from it: this is \textbf{an advantage of single-stage detectors}.\par}\vskip 4pt
\end{frame}

\framekicker{Section 12.3}
\begin{frame}[shrink=25]{Why single-stage is more robust to novelty}
\begin{minipage}[t]{0.494\textwidth}
\begin{icardwarn}
\cardh{Two-stage}
Forced to classify \textbf{small patches in isolation} --- which is extra difficult when a small patch of pixels looks nothing like the training data.
\end{icardwarn}
\end{minipage}%
\hfill
\begin{minipage}[t]{0.494\textwidth}
\begin{icardgood}
\cardh{Single-stage}
Can draw on features from \textbf{the entire input}, so it is more robust to novel test-time inputs.
\end{icardgood}
\end{minipage}%
{\fontsize{9.4}{12.6}\selectfont\color{fg2}Which connects back to chapter 5: robustness comes from \hilite{having more context to interpolate from}, not from having seen this exact style before.\par}\vskip 4pt
\end{frame}

\framekicker{Summary}
\begin{frame}[shrink=25]{What has to survive this chapter}
\begin{isteps}
\item \textbf{Detection when you do not need pixels} --- it is cheaper to run and far cheaper to label than segmentation.
\item \textbf{Two-stage proposes then classifies}; \textbf{single-stage predicts boxes and labels jointly} and is faster.
\item \textbf{YOLO divides the image into a grid}, and each cell predicts a box, a confidence, and a class distribution.
\item \textbf{Confidence is trained against IoU}, not a 0/1 flag --- so it measures \textit{how good} the box is, not only whether something is there.
\item \textbf{Input resolution matters} for detection, far more than for classification.
\end{isteps}
\end{frame}

\framekicker{Summary}
\begin{frame}[shrink=25]{\ldots{}and the two things worth carrying further}
\begin{minipage}[t]{0.494\textwidth}
\begin{icardwarn}
\cardh{An honest negative result}
The from-scratch model \textbf{stays undertrained}, and the chapter says so. Training a real COCO detector takes serious compute --- which is the argument for pretrained models, stated with evidence.
\end{icardwarn}
\end{minipage}%
\hfill
\begin{minipage}[t]{0.494\textwidth}
\begin{icardgood}
\cardh{Robustness has a structural cause}
RetinaNet handles a painting it never saw because it reads the \textbf{whole image}, and uses features from \textbf{several depths}.
\end{icardgood}
\end{minipage}%
\vskip 4pt
\reslink{NOTEBOOK}{03\_pretrained\_retinanet.ipynb}{../../course-slides/notebooks/ch12/03\_pretrained\_retinanet.ipynb}
\reslink{NEXT}{Chapter 13 --- Timeseries forecasting}{../ch13/index.html}
\vskip 2pt
\end{frame}

\end{document}
